% GENERATED FILE. Edit canonical lessons, then run npm run generate:books.
% canonical-source-hash: fnv1a64:896a0aa4e85a9b36
% canonical-lessons: FR-C427-climatisation, FR-C427-chauffage, FR-C427-radiateur, FR-C427-taie, FR-C427-peignoir

\chapter{Le confort de la chambre}
\label{ch:fr-le-confort-de-la-chambre}
\hlchaptermodality{427}

\begin{chapteropening}
\textbf{By the end of this chapter:} \emph{I can talk about comfort in a hotel room: the air conditioning, the heating, the radiator, the pillowcase and the bathrobe.}

Complete the last lesson of chapter 427: I can talk about comfort in a hotel room: the air conditioning, the heating, the radiator, the pillowcase and the bathrobe.
\end{chapteropening}

\section[la climatisation]{la climatisation — the air conditioning}

\label{lesson:FR-C427-climatisation}



\begin{quote}
Before the new one: say the French for the stay, then the French for to unpack.
\end{quote}

\subsection*{You'll want to know: la climatisation}
\textbf{la climatisation} — \textquotedblleft{}the air conditioning\textquotedblright{}.

\begin{quote}
La climatisation fait trop de bruit la nuit.
\end{quote}

\begin{quote}
\emph{The air conditioning makes too much noise at night.}
\end{quote}

In speech, \emph{la clim}. \emph{Climatisé} — air-conditioned.

\begin{cousinweb}
\emph{climatisation} is built on \emph{climat} 'climate', from Latin \emph{clima}, a learned word from Greek \emph{klíma} 'slope, region' (Lewis \& Short). EtymDB lists it without a source.
\end{cousinweb}

\begin{grammarlens}[title={What you've built}]
One more word for telling, explaining and describing at B1.
\end{grammarlens}

\subsection*{Your turn}
\begin{itemize}
\raggedright
  \item \emph{Say it:} \emph{la climatisation}
  \item \emph{Say it:} \emph{la climatisation}, once more
  \item \emph{Say it:} a sentence of your own with \emph{la climatisation}
  \item \emph{Recall:} say the French for the stay, then the French for to unpack, then say \emph{la climatisation} again
\end{itemize}

\begin{tcolorbox}[breakable,colback=teal!4,colframe=teal!35!black,title={Before you move on}]
What does la climatisation mean? (\textquotedblleft{}The air conditioning\textquotedblright{}.) Say it once more.
\end{tcolorbox}
\section[le chauffage]{le chauffage — the heating}

\label{lesson:FR-C427-chauffage}



\begin{quote}
Before the new one: say the French for to unpack, then the French for the air conditioning.
\end{quote}

\subsection*{You'll want to know: le chauffage}
\textbf{le chauffage} — \textquotedblleft{}the heating\textquotedblright{}.

\begin{quote}
En hiver, le chauffage marche toute la journée.
\end{quote}

\begin{quote}
\emph{In winter the heating runs all day.}
\end{quote}

\emph{Le chauffage central} — central heating.

\begin{cousinweb}
\emph{chauffage} is built on \emph{chauffer} 'to heat' with \emph{-age}. \emph{Chauffer} is inherited from Latin \emph{calefacere} 'to make warm', from \emph{calēre} 'to be warm' (Lewis \& Short).
\end{cousinweb}

\begin{grammarlens}[title={What you've built}]
One more word for telling, explaining and describing at B1.
\end{grammarlens}

\subsection*{Your turn}
\begin{itemize}
\raggedright
  \item \emph{Say it:} \emph{le chauffage}
  \item \emph{Say it:} \emph{le chauffage}, once more
  \item \emph{Say it:} a sentence of your own with \emph{le chauffage}
  \item \emph{Recall:} say the French for to unpack, then the French for the air conditioning, then say \emph{le chauffage} again
\end{itemize}

\begin{tcolorbox}[breakable,colback=teal!4,colframe=teal!35!black,title={Before you move on}]
What does le chauffage mean? (\textquotedblleft{}The heating\textquotedblright{}.) Say it once more.
\end{tcolorbox}
\section[le radiateur]{le radiateur — the radiator}

\label{lesson:FR-C427-radiateur}



\begin{quote}
Before the new one: say the French for the air conditioning, then the French for the heating.
\end{quote}

\subsection*{You'll want to know: le radiateur}
\textbf{le radiateur} — \textquotedblleft{}the radiator\textquotedblright{}.

\begin{quote}
Elle a mis ses chaussettes sur le radiateur.
\end{quote}

\begin{quote}
\emph{She put her socks on the radiator.}
\end{quote}

Also the radiator of a car.

\begin{cousinweb}
\emph{radiateur} is a learned formation on Latin \emph{radiāre} 'to send out rays', from \emph{radius} 'ray' (EtymDB).
\end{cousinweb}

\begin{grammarlens}[title={What you've built}]
One more word for telling, explaining and describing at B1.
\end{grammarlens}

\subsection*{Your turn}
\begin{itemize}
\raggedright
  \item \emph{Say it:} \emph{le radiateur}
  \item \emph{Say it:} \emph{le radiateur}, once more
  \item \emph{Say it:} a sentence of your own with \emph{le radiateur}
  \item \emph{Recall:} say the French for the air conditioning, then the French for the heating, then say \emph{le radiateur} again
\end{itemize}

\begin{tcolorbox}[breakable,colback=teal!4,colframe=teal!35!black,title={Before you move on}]
What does le radiateur mean? (\textquotedblleft{}The radiator\textquotedblright{}.) Say it once more.
\end{tcolorbox}
\section[la taie]{la taie — the pillowcase}

\label{lesson:FR-C427-taie}



\begin{quote}
Before the new one: say the French for the heating, then the French for the radiator.
\end{quote}

\subsection*{You'll want to know: la taie}
\textbf{la taie} — \textquotedblleft{}the pillowcase\textquotedblright{}.

\begin{quote}
La taie d'oreiller sent la lessive.
\end{quote}

\begin{quote}
\emph{The pillowcase smells of washing powder.}
\end{quote}

In full: \emph{la taie d'oreiller}.

\begin{cousinweb}
\emph{taie} comes from Old French \emph{teie} (EtymDB), from Latin \emph{thēca} 'case, cover', which came from Greek \emph{th\'{\={e}}kē} 'box' (Lewis \& Short).
\end{cousinweb}

\begin{grammarlens}[title={What you've built}]
One more word for telling, explaining and describing at B1.
\end{grammarlens}

\subsection*{Your turn}
\begin{itemize}
\raggedright
  \item \emph{Say it:} \emph{la taie}
  \item \emph{Say it:} \emph{la taie}, once more
  \item \emph{Say it:} a sentence of your own with \emph{la taie}
  \item \emph{Recall:} say the French for the heating, then the French for the radiator, then say \emph{la taie} again
\end{itemize}

\begin{tcolorbox}[breakable,colback=teal!4,colframe=teal!35!black,title={Before you move on}]
What does la taie mean? (\textquotedblleft{}The pillowcase\textquotedblright{}.) Say it once more.
\end{tcolorbox}
\section[le peignoir]{le peignoir — the bathrobe, the dressing gown}

\label{lesson:FR-C427-peignoir}



\begin{quote}
Before the new one: say the French for the radiator, then the French for the pillowcase.
\end{quote}

\subsection*{You'll want to know: le peignoir}
\textbf{le peignoir} — \textquotedblleft{}the bathrobe, the dressing gown\textquotedblright{}.

\begin{quote}
Après la douche, il met le peignoir de l'hôtel.
\end{quote}

\begin{quote}
\emph{After the shower he puts on the hotel bathrobe.}
\end{quote}

\emph{Un peignoir de bain} — a bathrobe.

\begin{cousinweb}
\emph{peignoir} is built on \emph{peigner} 'to comb', from Latin \emph{pectināre}, from \emph{pecten} 'comb' (Lewis \& Short): first the gown worn while your hair was combed.
\end{cousinweb}

\begin{grammarlens}[title={What you've built}]
One more word for telling, explaining and describing at B1.
\end{grammarlens}

\subsection*{Your turn}
\begin{itemize}
\raggedright
  \item \emph{Say it:} \emph{le peignoir}
  \item \emph{Say it:} \emph{le peignoir}, once more
  \item \emph{Say it:} a sentence of your own with \emph{le peignoir}
  \item \emph{Recall:} say the French for the radiator, then the French for the pillowcase, then say \emph{le peignoir} again
\end{itemize}

\begin{tcolorbox}[breakable,colback=teal!4,colframe=teal!35!black,title={Before you move on}]
What does le peignoir mean? (\textquotedblleft{}The bathrobe, the dressing gown\textquotedblright{}.) Say it once more.
\end{tcolorbox}
